% GENERATED FILE. Edit canonical lessons, then run npm run generate:books.
% canonical-source-hash: fnv1a64:4d1c742aa0255983
% canonical-lessons: GU-C552-manc, GU-C552-abhinay, GU-C552-darsak, GU-C552-gajavvu, GU-C552-jaamvu

\chapter{If the Show Goes On}
\label{ch:gu-if-the-show-goes-on}
\hlchaptermodality{552}

\begin{chapteropening}
\textbf{By the end of this chapter:} \emph{I can talk about a show, the stage and the audience.}

Complete the last lesson of chapter 552: I can talk about a show, the stage and the audience.
\end{chapteropening}

\section[\texorpdfstring{\gu{મંચ} (manc)}{મંચ (manc)}]{\gu{મંચ} (manc) — a stage}

\label{lesson:GU-C552-manc}



\begin{quote}
Before the new one: say the Gujarati for to make (someone) run, then the Gujarati for a trophy.
\end{quote}

\subsection*{You'll want to know: \gu{મંચ}}
\textbf{\gu{મંચ}} — \emph{manc} — \textquotedblleft{}a stage\textquotedblright{}.

\begin{quote}
\gu{જો} \gu{મંચ} \gu{તૈયાર} \gu{હશે} \gu{તો} \gu{નાટક} \gu{સમયસર} \gu{શરૂ} \gu{થશે}.
\end{quote}

\begin{quote}
\emph{jo manc taiyār haśe to nāṭak samaysar śarū thaśe.}
\end{quote}

\begin{quote}
If the stage is ready, the play will start on time.
\end{quote}

A masculine noun.

\begin{cousinweb}
\textbf{\gu{મંચ}} (\emph{manc}) is a tatsama, taken from Sanskrit \emph{mañca}, \textquotedblleft{}a raised platform, a stage\textquotedblright{} (Monier-Williams).
\end{cousinweb}

\begin{grammarlens}[title={What you've built}]
One more word for saying what happens if.
\end{grammarlens}

\subsection*{Your turn}
\begin{itemize}
\raggedright
  \item \emph{Say it:} \emph{manc}
  \item \emph{Say it:} \emph{manc}, once more
  \item \emph{Say it:} \emph{jo manc taiyār haśe to nāṭak samaysar śarū thaśe.} — If the stage is ready, the play will start on time.
  \item \emph{Recall:} say the Gujarati for to make (someone) run, then the Gujarati for a trophy, then say \emph{manc} again
\end{itemize}

\begin{tcolorbox}[breakable,colback=teal!4,colframe=teal!35!black,title={Before you move on}]
What does \gu{મંચ} mean? (\textquotedblleft{}A stage\textquotedblright{}.) Say it once more.
\end{tcolorbox}
\section[\texorpdfstring{\gu{અભિનય} (abhinay)}{અભિનય (abhinay)}]{\gu{અભિનય} (abhinay) — acting}

\label{lesson:GU-C552-abhinay}



\begin{quote}
Before the new one: say the Gujarati for a trophy, then the Gujarati for a stage.
\end{quote}

\subsection*{You'll want to know: \gu{અભિનય}}
\textbf{\gu{અભિનય}} — \emph{abhinay} — \textquotedblleft{}acting\textquotedblright{}.

\begin{quote}
\gu{જો} \gu{તારો} \gu{અભિનય} \gu{સારો} \gu{હશે} \gu{તો} \gu{તું} \gu{ફિલ્મોમાં} \gu{કામ} \gu{કરી} \gu{શકીશ}.
\end{quote}

\begin{quote}
\emph{jo tāro abhinay sāro haśe to tũ philmomā̃ kām karī śakīś.}
\end{quote}

\begin{quote}
If your acting is good, you'll be able to work in films.
\end{quote}

A masculine noun.

\begin{cousinweb}
\textbf{\gu{અભિનય}} (\emph{abhinay}) is a tatsama, taken from Sanskrit \emph{abhinaya}, the showing of a feeling by look and gesture: acting (Monier-Williams).
\end{cousinweb}

\begin{grammarlens}[title={What you've built}]
One more word for saying what happens if.
\end{grammarlens}

\subsection*{Your turn}
\begin{itemize}
\raggedright
  \item \emph{Say it:} \emph{abhinay}
  \item \emph{Say it:} \emph{abhinay}, once more
  \item \emph{Say it:} \emph{jo tāro abhinay sāro haśe to tũ philmomā̃ kām karī śakīś.} — If your acting is good, you'll be able to work in films.
  \item \emph{Recall:} say the Gujarati for a trophy, then the Gujarati for a stage, then say \emph{abhinay} again
\end{itemize}

\begin{tcolorbox}[breakable,colback=teal!4,colframe=teal!35!black,title={Before you move on}]
What does \gu{અભિનય} mean? (\textquotedblleft{}Acting\textquotedblright{}.) Say it once more.
\end{tcolorbox}
\section[\texorpdfstring{\gu{દર્શક} (darśak)}{દર્શક (darśak)}]{\gu{દર્શક} (darśak) — a spectator}

\label{lesson:GU-C552-darsak}



\begin{quote}
Before the new one: say the Gujarati for a stage, then the Gujarati for acting.
\end{quote}

\subsection*{You'll want to know: \gu{દર્શક}}
\textbf{\gu{દર્શક}} — \emph{darśak} — \textquotedblleft{}a spectator\textquotedblright{}.

\begin{quote}
\gu{જો} \gu{દર્શકો} \gu{ઓછા} \gu{હશે} \gu{તો} \gu{નાટક} \gu{બંધ} \gu{થઈ} \gu{જશે}.
\end{quote}

\begin{quote}
\emph{jo darśako ochā haśe to nāṭak bandh thaī jaśe.}
\end{quote}

\begin{quote}
If there are few spectators, the play will be stopped.
\end{quote}

A masculine or feminine noun, by the person meant.

\begin{cousinweb}
\textbf{\gu{દર્શક}} (\emph{darśak}) is a tatsama, taken from Sanskrit \emph{darśaka}, \textquotedblleft{}one who sees\textquotedblright{} (Monier-Williams).
\end{cousinweb}

\begin{grammarlens}[title={What you've built}]
One more word for saying what happens if.
\end{grammarlens}

\subsection*{Your turn}
\begin{itemize}
\raggedright
  \item \emph{Say it:} \emph{darśak}
  \item \emph{Say it:} \emph{darśak}, once more
  \item \emph{Say it:} \emph{jo darśako ochā haśe to nāṭak bandh thaī jaśe.} — If there are few spectators, the play will be stopped.
  \item \emph{Recall:} say the Gujarati for a stage, then the Gujarati for acting, then say \emph{darśak} again
\end{itemize}

\begin{tcolorbox}[breakable,colback=teal!4,colframe=teal!35!black,title={Before you move on}]
What does \gu{દર્શક} mean? (\textquotedblleft{}A spectator\textquotedblright{}.) Say it once more.
\end{tcolorbox}
\section[\texorpdfstring{\gu{ગજાવવું} (gajāvvũ)}{ગજાવવું (gajāvvũ)}]{\gu{ગજાવવું} (gajāvvũ) — to make ring (with noise)}

\label{lesson:GU-C552-gajavvu}



\begin{quote}
Before the new one: say the Gujarati for acting, then the Gujarati for a spectator.
\end{quote}

\subsection*{You'll want to know: \gu{ગજાવવું}}
\textbf{\gu{ગજાવવું}} — \emph{gajāvvũ} — \textquotedblleft{}to make ring (with noise)\textquotedblright{}.

\begin{quote}
\gu{જો} \gu{આપણી} \gu{ટીમ} \gu{જીતશે} \gu{તો} \gu{આપણે} \gu{આખું} \gu{મેદાન} \gu{ગજાવીશું}.
\end{quote}

\begin{quote}
\emph{jo āpṇī ṭīm jītśe to āpaṇe ākhũ medān gajāvīśũ.}
\end{quote}

\begin{quote}
If our team wins, we'll make the whole ground ring.
\end{quote}

Past, with \gu{તેણે} (he or she …): \gu{તેણે} \gu{ઓરડો} \gu{ગજાવ્યો}, \emph{teṇe orḍo gajāvyo}.

\begin{cousinweb}
\textbf{\gu{ગજાવવું}} (\emph{gajāvvũ}) is built from \textbf{\gu{ગાજવું}}, \textquotedblleft{}to thunder\textquotedblright{}, with the causative \emph{-āv-} and a shortened vowel: to make a place thunder. CDIAL files \textbf{\gu{ગાજવું}} under Sanskrit \emph{garjati}, \textquotedblleft{}roars\textquotedblright{} (CDIAL 4046).
\end{cousinweb}

\begin{grammarlens}[title={What you've built}]
One more word for saying what happens if.
\end{grammarlens}

\subsection*{Your turn}
\begin{itemize}
\raggedright
  \item \emph{Say it:} \emph{gajāvvũ}
  \item \emph{Say it:} \emph{gajāvvũ}, once more
  \item \emph{Say it:} \emph{jo āpṇī ṭīm jītśe to āpaṇe ākhũ medān gajāvīśũ.} — If our team wins, we'll make the whole ground ring.
  \item \emph{Recall:} say the Gujarati for acting, then the Gujarati for a spectator, then say \emph{gajāvvũ} again
\end{itemize}

\begin{tcolorbox}[breakable,colback=teal!4,colframe=teal!35!black,title={Before you move on}]
What does \gu{ગજાવવું} mean? (\textquotedblleft{}To make ring (with noise)\textquotedblright{}.) Say it once more.
\end{tcolorbox}
\section[\texorpdfstring{\gu{જામવું} (jāmvũ)}{જામવું (jāmvũ)}]{\gu{જામવું} (jāmvũ) — to get going (of a show or a party)}

\label{lesson:GU-C552-jaamvu}



\begin{quote}
Before the new one: say the Gujarati for a spectator, then the Gujarati for to make ring.
\end{quote}

\subsection*{You'll want to know: \gu{જામવું}}
\textbf{\gu{જામવું}} — \emph{jāmvũ} — \textquotedblleft{}to get going (of a show or a party)\textquotedblright{}.

\begin{quote}
\gu{જો} \gu{સંગીત} \gu{સારું} \gu{હશે} \gu{તો} \gu{કાર્યક્રમ} \gu{જામશે}.
\end{quote}

\begin{quote}
\emph{jo sangīt sārũ haśe to kāryakram jāmśe.}
\end{quote}

\begin{quote}
If the music is good, the show will really get going.
\end{quote}

Past (it …): \gu{તે} \gu{કાર્યક્રમ} \gu{જામી} \gu{ગયો}, \emph{te kāryakram jāmī gayo}.

\begin{cousinweb}
\textbf{\gu{જામવું}} (\emph{jāmvũ}) is filed by CDIAL under Sanskrit \emph{yamyátē}, \textquotedblleft{}is fixed\textquotedblright{}, and CDIAL lists the Gujarati verb, \textquotedblleft{}to collect, to set\textquotedblright{} (CDIAL 10428): milk sets into curd, and a show \textquotedblleft{}sets\textquotedblright{} when it takes hold.
\end{cousinweb}

\begin{grammarlens}[title={What you've built}]
One more word for saying what happens if.
\end{grammarlens}

\subsection*{Your turn}
\begin{itemize}
\raggedright
  \item \emph{Say it:} \emph{jāmvũ}
  \item \emph{Say it:} \emph{jāmvũ}, once more
  \item \emph{Say it:} \emph{jo sangīt sārũ haśe to kāryakram jāmśe.} — If the music is good, the show will really get going.
  \item \emph{Recall:} say the Gujarati for a spectator, then the Gujarati for to make ring, then say \emph{jāmvũ} again
\end{itemize}

\begin{tcolorbox}[breakable,colback=teal!4,colframe=teal!35!black,title={Before you move on}]
What does \gu{જામવું} mean? (\textquotedblleft{}To get going (of a show or a party)\textquotedblright{}.) Say it once more.
\end{tcolorbox}
